\section{Fórmula para $F_2$ e identificación del tamaño físico}
Para $b_0=v-1$ formemos los núcleos $G_v,T_v,U,H$ según el apartado 3 por convolución
normalizada (en este párrafo los subíndices de lado son $L,R$). Sean $j$ el núcleo diagonal
de $T-G$, $j_0=\tr(T_R-G_R)$ y $\sigma_3$ ahora la integral del cubo por entradas.
Definimos exactamente
\begin{align*}
 \phi(v)={}&j_0^2/2+\tr T_R^2-3j_0/2-3\langle j,j\rangle/4+\langle j,Hj\rangle/2\\
 &-(\sigma_3(T_L)+\sigma_3(T_R))/3+2\int b_0U^2+2\sigma_3(U)/3.
\end{align*}
Cada lado tiene masa uno y la unión disjunta masa dos. En el tamaño físico $k\Xi(P)=\phi(v)$:
cada entrada de matriz ordinaria de un núcleo vale $s$ veces el núcleo y cada suma
ordinaria sobre un lado aporta un factor $1/s$. Esto incluye todas las diagonales y las trazas.
Sea $Z$ el núcleo $z-1$, y $O$ el bloque simétrico de diagonal nula con bloque cruzado $Z$.
Para $y$ cerca de uno escribimos
\[
 d(y)=\det(I-y^2 Z^\top Z)^{-1/2},\qquad H_y=yO(I+yO)^{-1}.
\]
El propagador $H_y$ es un núcleo acotado: se escribe $yO-y^2 O(I+yO)^{-1}O$;
no contiene un núcleo identidad. La brecha da un disco de $y$ de radio fijo.
Para prolongar $V_j$, $j=1,2$, de ronda 2, parte 10, sumamos sobre $m$ vértices
distintos, grados $d_i\ge3$, $2E=\sum d_i$, $E-m=j$. Ordenando los vértices,
el peso es $\prod_i[w^{d_i}](1+w)e^{-w+w^2/2}$, dividido por $m!$, multiplicado
por la suma sobre todos los emparejamientos de las semiaristas etiquetadas.
Se incluyen los desconectados. En cada contracción se impone que los vértices sean
distintos por inversión sobre particiones $\pi$ de los vértices:
\[
 \sum_\pi\left[\prod_{C\in\pi}(-1)^{|C|-1}(|C|-1)!\right]
 s^{E-|\pi|}\int\prod_{uv\ {\rm aristas}} H_y(x_{\pi(u)},x_{\pi(v)})\,d\boldsymbol x.
\]
Se integra una vez por bloque, sobre la unión disjunta. Los vértices de lados distintos
nunca se identifican al fijar sus lados. En el tamaño físico ésta es exactamente la
suma de índices distintos; tras replicar, la distinción es entre copias individuales.
Hay $m\le2j$, $E\le3j$ y $E-|\pi|\ge j$. La fórmula define $V_j(y,s)$ con coeficientes
acotados e invariantes al replicar.
Definamos los momentos centrales de la mezcla y su aproximación
\begin{align*}
 \mu_l(s)&=\sum_{q=0}^l(-1)^{l-q}\binom lq\prod_{a=0}^{q-1}(1-as)^{-1},\\
 H_*(s)&=\sum_{l=0}^5\frac{\mu_l(s)}{l!}
 \left.\partial_y^l[d(y)(1+V_1(y,s)+V_2(y,s))]\right|_{y=1}.
\end{align*}
Para $s=1/k$ son los de la ronda 2 ($k$ suficientemente grande). Sea
$c(x)=\tfrac12\log(1-2x)-2\log(1-x)$. El coeficiente buscado es
$F_2(z)=[s^2]\Gamma_z(s)$, con
\begin{align*}
 \Gamma_z(s)={}&s/6+\log H_*(s)
 +(\langle\alpha\rangle+\langle\beta\rangle)/s
 +(\langle c(sv)\rangle-s)/s^2\\
 &+(\langle\log a\rangle+\langle\log b\rangle)/(2s)
 +\tfrac12\log\det(I-W^\top W)-\tfrac12\log\bar\ell-s\phi(v).
\end{align*}
Para verificar signos: el término $-b$ de ronda 8, parte 02 (potenciales del
escalado exacto $D$ a partir de $P(1-P)$), vale
$(\langle\alpha\rangle+\langle\beta\rangle)/s$. En ese apartado se reemplaza la
aproximación de $\log(k^k\per D/k!)$ por $\log H_*$, con error uniforme $O(s^3)$
según ronda 2, parte 10. Stirling, el factor $\sqrt{4\pi k}$ y $\log2/2$ dejan
$-k+1/(12k)+O(k^{-3})$. En $R=g_T+1/(12k)-\Xi$ esto da $s/6$ y $-s/s^2$.
Así $R=\Gamma_z(1/k)+O(k^{-3})$ uniformemente.
Los cocientes en cero son removibles: $c(x)=x+O(x^3)$ y los márgenes de $v$ valen uno.
Los determinantes en cero están separados de cero por la brecha y
$\tr(Z^\top Z)=O(1)$; $H_*(0)=d(1)$. Tomamos las ramas reales positivas en el centro.
La acotación holomorfa en un disco fijo se detalla en el apartado siguiente.
